\chapter{Quaternions.}

\section{What are they exactly?}

Quaternions are closely related to complex numbers. However, whereas complex numbers are comprised of one real component and one
imaginary component, quaternions are comprised of one real component and three imaginary components, for a total of four
components altogether. \\

Quaternions can be thought of as extending the complex number system from two dimensions into four dimensions. Whereas
complex numbers reside in the 2-dimensional space which is denoted as \(\mathbb{C}\) (which stands for ``complex''),
quaternions reside in a 4-dimensional number space which is denoted as \(\mathbb{H}\) (which stands for ``Hamilton''). \\

Since quaternions are comprised of multiple imaginary components, they are categorised as belonging to a class of
numbers known as ``hypercomplex numbers''. Furthermore, since they are comprised of four components in total, this helps
to explain why they are called ``quaternions''; as ``quat'' means of or relating to four parts. \(\mathbb{C}\)

A quaternion is often denoted using a form which is similar to the following;

\begin{align*}
Q & = w   + \qi x   + \qi y   + \quaternion{k} z  \\
Q & = a   + \qi b   + \qj c   + \quaternion{k} d  \\
Q & = Q_0 + \qi Q_1 + \qj Q_2 + \quaternion{k} Q_3
\end{align*}
\vspace{0.0\baselineskip}

The first component, or the real component of the quaternion, i.e. \(w\), \(a\), or \(Q_0\) in the previous examples,
is often referred to as the scalar part of the quaternion, while the remaining three components are collectively
referred to as the vector part of the quaternion.

\begin{equation*}
Q = \underbrace{w}_{\text{scalar}} + \underbrace{\qi x + \qj y + \quaternion{k}z}_{\text{vector}}
\end{equation*}

If we were to denote the vector component of the quaternion \(Q\) as \(\vec{Q}\), then we could rewrite the three example
quaternions from above as follows;

\begin{align*}
    Q & = w + \vec{Q}  \\
    Q & = a + \vec{Q}  \\
    Q & = Q_0 + \vec{Q}
\end{align*}
\vspace{0.0\baselineskip}

where;

\begin{align*}
    \vec{Q} &= (\qi x, \qj y, \quaternion{k}z) \; \text{or}  \\
            &= (\qi b, \qj c, \quaternion{k}d) \; \text{or}  \\
            &= (\qi Q_{1}, \qj Q_{2}, \quaternion{k}Q_{3})
\end{align*}


\subsection{Quaternion basis units and the rules governing their multiplication.}

Let's recall our three example quaternions from earlier, that is;

\begin{align*}
Q & = w   + \qi x   + \qi y   + \quaternion{k} z  \\
Q & = a   + \qi b   + \qj c   + \quaternion{k} d  \\
Q & = Q_0 + \qi Q_1 + \qj Q_2 + \quaternion{k} Q_3
\end{align*}
\vspace{0.0\baselineskip}

In these three examples, \(1\), \(\qi\), \(\qj\), and \(\quaternion{k}\) are all known as the quaternion basis units,
even though only \(\qi\), \(\qj\), and \(\quaternion{k}\) are explicitly depicted in the three definitions of \(Q\). \\

When it comes to multiplying these units with themselves and each other, they satisfy the following rules;

\begin{equation*}
\qi^2=\qj^2=\quaternion{k}^2=\quaternion{ijk}=-1
\end{equation*}
\vspace{0.0\baselineskip}

From this, the following additional multiplication rules can also be derived;

\begin{equation*}
\quaternion{ij}=\quaternion{k},\quad \quaternion{jk}=\qi,\quad \quaternion{ki}=\qj,
\end{equation*}

as well as;

\begin{equation*}
\quaternion{ji}=-\quaternion{k}, \quad \quaternion{kj}=-\qi,\quad \quaternion{ik}=-\qj
\end{equation*}


\subsection{Pure quaternions.}

A pure quaternion -- sometimes also referred to as a pure imaginary quaternion or a vector quaternion, is a quaternion
whose real or scalar component is equal to zero. That is, if we have;

\begin{equation*}
    Q = Q_0 + Q_1 + Q_2 + Q_3 \; \; \; \; \; \; \mid \; \; Q_0 = 0
\end{equation*}
\vspace{0.0\baselineskip}

then;

\begin{align*}
    Q & = 0 + Q_1 + Q_2 + Q_3  \\
      & = \vec{Q}
\end{align*}
\vspace{0.0\baselineskip}


\subsection{Unit quaternions.}

A unit quaternion -- not to be confused with a quaternion unit such as \(1\), \(\qi\), \(\qj\), or \(\qk\), is a
quaternion whose length or magnitude is equal to one. For example, the following quaternions are all unit quaternions,
because they all have a length or magnitude of one;

\begin{align*}
    Q1 & = 1  \\
    Q2 & = 0.5 + \qi 0.5 + \qj 0.5 + \quaternion{k}0.5  \\
    Q3 & = 0.707 + \qi 0.707  \\
    Q4 & = \qi 0.577 + \qj 0.577 + \quaternion{k}0.577
\end{align*}
\vspace{0.0\baselineskip}

To illustrate, first consider the quaternion \(Q2\).

\begin{align*}
    \lVert Q2 \rVert & = \sqrt{0.5^2 + 0.5^2 + 0.5^2 + 0.5^2}  \\
                     & = \sqrt{0.25 + 0.25 + 0.25 + 0.25}  \\
                     & = \sqrt{1}  \\
                     & = 1
\end{align*}
\vspace{0.0\baselineskip}

Next consider the quaternion \(Q4\).

\begin{align*}
    \lVert Q4 \rVert & = \sqrt{0.577^2 + 0.577^2 + 0.577^2}  \\
                     & = \sqrt{0.333 + 0.333 + 0.333}  \\
                     & = \sqrt{1}  \\
                     & = 1
\end{align*}
\vspace{0.0\baselineskip}

It might be worth noting - but then again it might not, that when taken altogether, all of the unit quaternions form the
surface of a 4-dimensional hypersphere. I mention that this might not be worth noting, because the concept of a
4-dimensional hypersphere seems almost incomprehensible to some people nd they would rather not have to think about it!
